
\documentclass{article}
\usepackage{colortbl}
\usepackage{makecell}
\usepackage{multirow}
\usepackage{supertabular}

\begin{document}

\newcounter{utterance}

\centering \large Interaction Transcript for game `guesswhat', experiment `Level\_3', episode 7 with Qwen3{-}32B{-}t0.0.
\vspace{24pt}

{ \footnotesize  \setcounter{utterance}{1}
\setlength{\tabcolsep}{0pt}
\begin{supertabular}{c@{$\;$}|p{.15\linewidth}@{}p{.15\linewidth}p{.15\linewidth}p{.15\linewidth}p{.15\linewidth}p{.15\linewidth}}
    \# & $\;$A & \multicolumn{4}{c}{Game Master} & $\;\:$B\\
    \hline

    \theutterance \stepcounter{utterance}  
    & & \multicolumn{4}{p{0.6\linewidth}}{
        \cellcolor[rgb]{0.9,0.9,0.9}{
            \makecell[{{p{\linewidth}}}]{
                \texttt{\tiny{[P1$\langle$GM]}}
                \texttt{Let's play a guessing game. Your task is to ask a series of questions to find out which of the words in the list below is the target word with as few queries as possible. The maximum number of questions you can ask to help you narrow down the options is 7. Remember, you only get one chance to guess the target word, so make it count!} \\
\\ 
\texttt{The rules of the game are: } \\
\\ 
\texttt{1. Only ask one question per turn. DO NOT ASK MORE THAN ONE QUESTION.} \\
\\ 
\texttt{2. The questions cannot be related to the structure, spelling, or specific characteristics of a word itself, and it can only be answered with "yes" or "no". } \\
\\ 
\texttt{3. Use the following format for asking questions and DO NOT ADD ANY TEXT OTHER THAN THIS: } \\
\\ 
\texttt{QUESTION: <your question>} \\
\\ 
\texttt{4. You have only one opportunity to guess the target word. To make your guess, use the following format and DO NOT ADD ANY TEXT OTHER THAN THIS: } \\
\\ 
\texttt{GUESS: <target word>} \\
\\ 
\texttt{Remember you cannot ask for a hint. Make your question brief. Let’s start!} \\
\\ 
\texttt{List of candidate words: {[}'Canoe', 'Dinghy', 'Battleship', 'Destroyer', 'Tugboat', 'Ferry', 'Pontoon', 'Houseboat'{]}} \\
            }
        }
    }
    & \\ \\

    \theutterance \stepcounter{utterance}  
    & \multicolumn{4}{p{0.6\linewidth}}{
        \cellcolor[rgb]{0.9,0.9,0.9}{
            \makecell[{{p{\linewidth}}}]{
                \texttt{\tiny{[P1$\rangle$GM]}}
                \texttt{<think>} \\
\texttt{Okay, let's see. I need to figure out which word is the target by asking yes/no questions. The list includes Canoe, Dinghy, Battleship, Destroyer, Tugboat, Ferry, Pontoon, Houseboat. All of them are types of boats or ships, but they have different purposes and sizes.} \\
\\ 
\texttt{First, I should try to categorize them. Maybe start by asking if the target is a military vessel. Battleship and Destroyer are definitely military. The others are more civilian or non{-}military. So if I ask, "Is the target a military vessel?" and the answer is yes, then it's either Battleship or Destroyer. If no, then it's one of the others. That could split the list into two groups. Let me check: military vs non{-}military. That's a good first question. If the answer is yes, then next question could be about size or role. If no, then maybe ask about passenger capacity or function. For example, if it's not military, maybe ask if it's used for transportation of people (Ferry, Houseboat) versus other purposes like Tugboat, Canoe, etc. Let's start with the first question.} \\
\texttt{</think>} \\
\\ 
\texttt{QUESTION: Is the target a military vessel?} \\
            }
        }
    }
    & & \\ \\

    \theutterance \stepcounter{utterance}  
    & & & \multicolumn{2}{p{0.3\linewidth}}{
        \cellcolor[rgb]{0.9,0.9,0.9}{
            \makecell[{{p{\linewidth}}}]{
                \texttt{\tiny{[GM$|$GM]}}
                \texttt{Invalid format. Guesser must use the form 'QUESTION: ' or 'GUESS: '.} \\
            }
        }
    }
    & & \\ \\

    \theutterance \stepcounter{utterance}  
    & & & \multicolumn{2}{p{0.3\linewidth}}{
        \cellcolor[rgb]{0.9,0.9,0.9}{
            \makecell[{{p{\linewidth}}}]{
                \texttt{\tiny{[GM$|$GM]}}
                \texttt{abort game} \\
            }
        }
    }
    & & \\ \\

\end{supertabular}
}

\end{document}
